\uuid{6pL1}
\titre{Production industrielle : loi normale }

\niveau{}
\module{ Probabilités }
\chapitre{Lois usuelles}
\sousChapitre{Lois normales}
\theme{}
\auteur{Quentin Liard}
\datecreate{ 2026-04-06}
\organisation{AMSCC}
\difficulte{3}

\contenu{

\texte{ 
Le diamètre d’une pièce produite par une machine est modélisé par une variable aléatoire $X$ suivant une loi normale de moyenne $50$ mm et d’écart-type $0{,}4$ mm :
$$
X\sim \mathcal N(50\,;\,0{,}4).
$$

Une pièce est déclarée conforme si son diamètre appartient à l’intervalle $[49{,}5\,;\,50{,}5]$.
}

\begin{enumerate}

\item   \question{Calculer la probabilité qu’une pièce produite soit conforme.}
\indication{}
\reponse{
Une pièce est conforme si
\[
49{,}5 \le X \le 50{,}5.
\]
On centre et on réduit :
\[
Z=\frac{X-50}{0{,}4}\sim \mathcal N(0,1).
\]
Ainsi
\[
P(49{,}5\le X\le 50{,}5)
=
P\left(\frac{49{,}5-50}{0{,}4}\le Z\le \frac{50{,}5-50}{0{,}4}\right)
=
P(-1{,}25\le Z\le 1{,}25).
\]
Donc
\[
P(49{,}5\le X\le 50{,}5)=2\Phi(1{,}25)-1.
\]
Avec la table de la loi normale centrée réduite,
\[
\Phi(1{,}25)\approx 0{,}8944.
\]
Donc
\[
P(49{,}5\le X\le 50{,}5)\approx 2\times 0{,}8944-1=0{,}7888.
\]
La probabilité qu’une pièce soit conforme vaut donc environ
\[
0{,}789.
\]
}

\item   \question{Calculer la probabilité qu’une pièce soit trop petite.}
\indication{}
\reponse{
Une pièce est trop petite si
\[
X<49{,}5.
\]
Donc
\[
P(X<49{,}5)=P\left(Z<\frac{49{,}5-50}{0{,}4}\right)=P(Z<-1{,}25).
\]
Par symétrie de la loi normale,
\[
P(Z<-1{,}25)=1-\Phi(1{,}25)\approx 1-0{,}8944=0{,}1056.
\]
Ainsi,
\[
P(X<49{,}5)\approx 0{,}106.
\]
}

%\item   \question{Calculer la probabilité qu’une pièce soit trop grande.}
%\indication{}
%\reponse{
%Une pièce est trop grande si
%\[
%X>50{,}5.
%\]
%Donc
%\[
%P(X>50{,}5)=P\left(Z>\frac{50{,}5-50}{0{,}4}\right)=P(Z>1{,}25).
%\]
%Or
%\[
%P(Z>1{,}25)=1-\Phi(1{,}25)\approx 0{,}1056.
%\]
%Donc
%\[
%P(X>50{,}5)\approx 0{,}106.
%\]
%}

\item   \question{Déterminer la probabilité que le diamètre d’une pièce s’écarte de plus de $0{,}8$ mm de la moyenne.}
\indication{}
\reponse{
On cherche
\[
P(|X-50|>0{,}8).
\]
Cela équivaut à
\[
P(X<49{,}2 \ \text{ou}\ X>50{,}8).
\]
En centrant et réduisant :
\[
P(|X-50|>0{,}8)=P\left(\left|Z\right|>\frac{0{,}8}{0{,}4}\right)=P(|Z|>2).
\]
Donc
\[
P(|X-50|>0{,}8)=2P(Z>2)=2(1-\Phi(2)).
\]
Or
\[
\Phi(2)\approx 0{,}9772.
\]
Ainsi
\[
P(|X-50|>0{,}8)\approx 2(1-0{,}9772)=0{,}0456.
\]
Donc la probabilité cherchée vaut environ
\[
0{,}046.
\]
}

\item   \question{On prélève une pièce au hasard. Sachant qu’elle n’est pas conforme, quelle est la probabilité qu’elle soit trop grande ?}
\indication{}
\reponse{
On cherche
\[
P(X>50{,}5 \mid X\notin [49{,}5;50{,}5]).
\]
Or l’événement ``non conforme'' est
\[
\{X<49{,}5\}\cup\{X>50{,}5\},
\]
et les deux événements ont la même probabilité par symétrie de la loi normale autour de la moyenne $50$.

Donc
\[
P(X>50{,}5 \mid X\notin [49{,}5;50{,}5])
=
\frac{P(X>50{,}5)}{P(X<49{,}5)+P(X>50{,}5)}
=\frac{0{,}1056}{0{,}1056+0{,}1056}
=\frac12.
\]
Ainsi, la probabilité cherchée vaut
\[
0{,}5.
\]
}

%\item   \question{On prélève indépendamment $10$ pièces. On note $Y$ le nombre de pièces conformes. Quelle est la loi de $Y$ ? Préciser son paramètre.}
%\indication{}
%\reponse{
%Chaque pièce a la probabilité
%\[
%p=P(\text{une pièce est conforme})\approx 0{,}7888.
%\]
%Les $10$ prélèvements étant indépendants, $Y$, nombre de pièces conformes parmi $10$, suit une loi binomiale de paramètres $10$ et $p$ %:
%\[
%Y\sim \mathcal B(10,\,0{,}7888).
%\]
%On peut aussi écrire, de façon exacte,
%\[
%Y\sim \mathcal B\bigl(10,\ 2\Phi(1{,}25)-1\bigr).
%\]
%}

\end{enumerate}

}